% GENERATED FILE. Edit canonical lessons, then run npm run generate:books.
% canonical-source-hash: fnv1a64:effc9d3dcafecae2
% canonical-lessons: AR-C28-dhahaba, AR-C28-jaa

\chapter{Going and Coming}
\label{ch:ar-going-and-coming}
\hlchaptermodality{28}

\begin{chapteropening}
\textbf{By the end of this chapter:} \emph{I can say the Arabic for 'he went' and 'he came', read the three-consonant root under each, and explain why Arabic dictionaries list a verb as 'he did' rather than 'to do'.}

Say the going/coming pair, name the weak letter that vanishes from jāʾa and point to where it comes back, and restate the root-and-pattern engine dhahaba opened.
\end{chapteropening}

\section[dhahaba]{\ar{ذهب} (dhahaba) — \textquotedblleft{}he went,\textquotedblright{} and the engine under every Arabic word}

\label{lesson:AR-C28-dhahaba}



\begin{quote}
Every word so far has been a thing, a quality, or a fixed courtesy. Now a \textbf{verb} — and the machine that builds nearly every word in the language.
\end{quote}

\subsection*{You'll want to know — The letters in this word}
Three letters, one of them new.

\begin{itemize}
\raggedright
  \item \textbf{\ar{د}} (\emph{dāl}) — a small hook resting on the line, sounding like English \emph{d}. It is a \textbf{non-joiner}: like \emph{alif}, it will not connect to whatever follows it.
  \item \textbf{\ar{ذ}} (\emph{dhāl}) — that same hook with \textbf{one dot above}. The dot is the entire difference. Its sound is the \emph{th} of English \textbf{this} — voiced — never the \emph{th} of \emph{thin}.
  \item \textbf{\ar{ه}} (\emph{hāʾ}) and \textbf{\ar{ب}} (\emph{bāʾ}) are already yours.
\end{itemize}

Right to left: \textbf{\ar{ذ}} stands alone, then \textbf{\ar{هب}} joins into one piece. \emph{Dha-ha-ba}.

\begin{cousinweb}
\textbf{\ar{ذَهَبَ}} (\emph{dhahaba}) = \textquotedblleft{}\textbf{he went}.\textquotedblright{} Its root is the bare consonant skeleton \textbf{\ar{ذ}-\ar{ه}-\ar{ب}}. Keep those three consonants, pour a different set of vowels between them, and out comes a different word carved from the same idea:

\begin{itemize}
\raggedright
  \item \textbf{\ar{ذَهَبَ}} (\emph{dhahaba}) — \textquotedblleft{}he went.\textquotedblright{}
  \item \textbf{\ar{ذَهَب}} (\emph{dhahab}) — \textquotedblleft{}\textbf{gold}.\textquotedblright{}
  \item \textbf{\ar{مَذْهَب}} (\emph{madhhab}) — \textquotedblleft{}a \textbf{way one goes},\textquotedblright{} and so the word for a \textbf{school of Islamic law}.
\end{itemize}

That last one is not luck. Putting \emph{ma-} in front of a root gives the \textbf{place or the way} of its action, and English already owns two words from that same place-noun family: \textbf{\ar{مَسْجِد}} (\emph{masjid}), \textquotedblleft{}place of prostration,\textquotedblright{} worn down through Spanish into \textbf{mosque}; and \textbf{\ar{مَغْرِب}} (\emph{maghrib}), \textquotedblleft{}place of sunset,\textquotedblright{} kept whole as the \textbf{Maghreb}. The pattern is already in your vocabulary. Only the root is new.
\end{cousinweb}

\begin{grammarlens}[title={Grammar Lens — Arabic has no \textquotedblleft{}to go\textquotedblright{}}]
English dictionaries list a verb as \textquotedblleft{}\textbf{to go}.\textquotedblright{} Arabic dictionaries cannot — the language has no infinitive. The form they list is \textbf{\ar{ذَهَبَ}}, \textquotedblleft{}\textbf{he went}\textquotedblright{}: third person, masculine, singular, finished. Grammarians name each vowel pattern with a dummy root, \textbf{\ar{ف}-\ar{ع}-\ar{ل}} (\emph{f-ʿ-l}, \textquotedblleft{}to do\textquotedblright{}), so this shape is called \textbf{\ar{فَعَلَ}} (\emph{faʿala}). Every verb in these three chapters wears it.
\end{grammarlens}

\subsection*{Your turn}
Say these aloud:

\begin{itemize}
\raggedright
  \item right to left — \ar{ذ}, then \ar{هب} — \textquotedblleft{}dhahaba,\textquotedblright{} he went
  \item the root — dh-h-b — then dhahab, \textquotedblleft{}gold,\textquotedblright{} and madhhab, \textquotedblleft{}a school of law\textquotedblright{}
  \item the citation form — Arabic lists \textquotedblleft{}he went,\textquotedblright{} never \textquotedblleft{}to go\textquotedblright{}
\end{itemize}

\begin{tcolorbox}[breakable,colback=teal!4,colframe=teal!35!black,title={Before you move on}]
What single mark separates \textbf{\ar{د}} from \textbf{\ar{ذ}}, and what does it do to the sound? (\textbf{One dot above} — it turns \emph{d} into the \emph{th} of \textquotedblleft{}this.\textquotedblright{}) Name two other words built on \textbf{\ar{ذ}-\ar{ه}-\ar{ب}}. (\emph{\textbf{Dhahab}}, \textquotedblleft{}gold,\textquotedblright{} and \emph{\textbf{madhhab}}, \textquotedblleft{}a school of law.\textquotedblright{}) In what form does an Arabic dictionary list a verb? (\textbf{\textquotedblleft{}He went\textquotedblright{}} — \emph{dhahaba} — because Arabic has no infinitive.)
\end{tcolorbox}
\section[jāʾa]{\ar{جاء} (jāʾa) — \textquotedblleft{}he came,\textquotedblright{} and a root that hides one of its own letters}

\label{lesson:AR-C28-jaa}



\begin{quote}
\emph{Dhahaba} was the tidy case: three consonants, three vowels, nothing hidden. Its opposite number is where the engine gets interesting.
\end{quote}

\subsection*{You'll want to know — The letters in this word}
\textbf{No new letters here} — this word is built entirely from shapes you already have.

\begin{itemize}
\raggedright
  \item \textbf{\ar{ج}} (\emph{jīm}) — the hook-and-tail body with \textbf{one dot below}, sounding like the \emph{j} of English \emph{jam}.
  \item \textbf{\ar{ا}} (\emph{alif}) — the clean downstroke, here carrying a long \emph{ā}.
  \item \textbf{\ar{ء}} (\emph{hamza}) — the glottal stop, the catch in \textquotedblleft{}uh-oh,\textquotedblright{} sitting on the line by itself at the end.
\end{itemize}

Right to left: \textbf{\ar{جا}} joined, then \textbf{\ar{ء}} standing apart. \emph{Jāʾa} — two beats, with a hard little stop at the end.

\begin{cousinweb}
\textbf{\ar{جاء}} (\emph{jāʾa}) = \textquotedblleft{}\textbf{he came}.\textquotedblright{} Its root is \textbf{\ar{ج}-\ar{ي}-\ar{أ}} — and one of those three consonants is missing from the word you just read. That is the point.

Arabic calls \textbf{\ar{و}} and \textbf{\ar{ي}} \textbf{weak} letters. Poured into the \emph{faʿala} shape, the root would give \emph{jayaʾa}, which no Arab mouth wants to say; the \textbf{\ar{ي}} dissolves and its vowel stretches into the long \textbf{ā} you hear. The skeleton is still there, though, and it walks straight back out in \textbf{\ar{مَجِيء}} (\emph{majīʾ}), \textquotedblleft{}\textbf{an arrival}\textquotedblright{} — the same \emph{ma-} place-and-way prefix met in \emph{madhhab}, with the \textbf{\ar{ي}} restored in plain sight.

So the irregularity is not chaos. It is a second rule, applied everywhere a weak letter lands in the middle. Being honest about one thing: \textbf{\ar{جاء}} is the written, formal verb. To tell someone \textquotedblleft{}come here\textquotedblright{} in ordinary speech you say \textbf{\ar{تَعالَ}} (\emph{taʿāla}), which comes from a completely \textbf{different root} meaning \textquotedblleft{}to be high\textquotedblright{} — originally \textquotedblleft{}come up here.\textquotedblright{}
\end{cousinweb}

\subsection*{Your turn}
Say these aloud:

\begin{itemize}
\raggedright
  \item the pair — \textquotedblleft{}dhahaba, jāʾa\textquotedblright{} — he went, he came
  \item the vanished letter — root j-y-ʾ, and the \ar{ي} back in place in \textquotedblleft{}majīʾ\textquotedblright{}
  \item the everyday imperative — \textquotedblleft{}taʿāla,\textquotedblright{} from another root entirely
\end{itemize}

\begin{tcolorbox}[breakable,colback=teal!4,colframe=teal!35!black,title={Before you move on}]
Which letter of \textbf{\ar{جاء}}'s root never appears in the word itself, and where does it come back? (\textbf{\ar{ي}} — restored in \textbf{\ar{مَجِيء}}, \emph{majīʾ}, \textquotedblleft{}an arrival.\textquotedblright{}) What are Arabic's two \textbf{weak} letters? (\textbf{\ar{و}} and \textbf{\ar{ي}}.) What word does an Arabic speaker actually use for \textquotedblleft{}come here,\textquotedblright{} and is it from this root? (\emph{\textbf{Taʿāla}} — \textbf{no}, a different root meaning \textquotedblleft{}to be high.\textquotedblright{})
\end{tcolorbox}
